\chapter{Split the pile}

Hot wallets maximize convenience and online exposure.
Cold wallets reduce remote attack surface and add
friction. Match the fund to the exposure. Do not keep
savings on the same hot path as everyday minting.

\section{Hot}

A hot wallet lives on a phone, a browser extension, or a
desktop that is on the network. Fine for spend. Fine for
learning. Wrong for the pile.

Download only from the official site or store listing
you verified twice. Open source helps. Independent
reviews help more if you are not going to read the code.
Wallet Scrutiny and similar scorecards are starting
points, not blessings.

\section{Cold}

A cold wallet keeps the private key off the internet
while it signs. Hardware wallets are the usual form.
Paper wallets and air-gapped software exist. Brain
wallets exist as a way to lose funds to a guessable
sentence.

Most hardware still talks to a computer, over USB,
Bluetooth, QR, or an SD card. That does not make it hot.
The test is whether the key stays isolated from the
infected laptop during signing.

Multisig and account abstraction are authorization
models, not types of cold wallet. Either can be hot,
cold, or mixed.

\section{Small balances}

If loss would sting and not wreck you, a hot wallet is
the honest setup. Unique seed. Offline backup the same
day you create it. Careful signing from day one.

A small balance is practice for a larger one. Do not
learn seed handling after the number gets interesting.

\section{When the number grows}

Move the bulk onto hardware. Keep a hot wallet for
spend. The savings wallet does not need to mint, vote,
or approve random contracts.

If one person losing a laptop would change a life, you
are past ``one hot wallet'' even if the number still
feels medium.
